%!TEX program = xelatex
\documentclass[lang=en,11pt,a4paper,citestyle =authoryear]{elegantpaper}

% 标题
\title{HW1 - STAT 709}
\author{Wells Guan}

% 本文档命令
\usepackage{array,url,stix}
\usepackage{tikz}
\usepackage{tikz-cd}
\usepackage{enumitem}
\usepackage{amsmath}
\usepackage{amssymb}
\usepackage{listings}
\usepackage{subfigure}

\lstset{
    frame=single,
    columns=fullflexible,
    keepspaces=true,
    showstringspaces=false,
    breaklines=true,
    breakatwhitespace=false,
    basicstyle=\ttfamily\footnotesize,
    numbers=left,
    numberstyle=\tiny,
    numbersep=6pt
}
\newcommand{\ccr}[1]{\makecell{{\color{#1}\rule{1cm}{1cm}}}}
\newcommand{\code}[1]{\lstinline{#1}}
\newcommand{\prvd}{$\hfill \qedsymbol$}
\newcommand{\Z}{\mathbb{Z}}
\newcommand{\R}{\mathbb{R}}
\newcommand{\N}{\mathbb{N}}
\newcommand{\C}{\mathbb{C}}
\newcommand{\Q}{\mathbb{Q}}
\newcommand{\M}{\mathcal{M}}
\newcommand{\B}{\mathcal{B}}
\newcommand{\X}{\mathcal{X}}
\newcommand{\Hil}{\mathcal{H}}
\newcommand{\range}{\mathcal{R}}
\newcommand{\nul}{\mathcal{N}}

% 文档区
\begin{document}

% 标题
\maketitle

\noindent

\subsection*{Problem 1. Ex. 1.7. } 

 Let $\nu_i, i = 1, 2, \cdots$, be measures on $(\Omega, \mathcal{F})$ and $a_i, i = 1, 2, \cdots$, be
positive numbers. Show that $a_1\nu_1 +a_2\nu_2 + \cdots$ is a measure on $(\Omega, \mathcal{F})$.\par
\vspace{0.5em}
\textbf{Sol.} \par
Denote $\sum\limits_{i=1}^{\infty}a_i\nu_i =: \nu$ and to prove $\nu(\varnothing) = 0$, we have
\[
\nu(\varnothing) = \sum\limits_{i=1}^{\infty}a_i\nu_i(\varnothing) = 0
\]
and for any disjoint $\{E_i\}_{i=1}^{\infty}\subset\mathcal{F}$, we have
\[
\begin{aligned}
    \nu\left(\bigsqcup_{i=1}^{\infty}E_i\right) &=\sum\limits_{j=1}^{\infty}a_j\nu_j\left(\bigsqcup_{i=1}^{\infty}E_i\right)\\
    &= \sum\limits_{i=1}^{\infty}\sum\limits_{j=1}^{\infty} a_j\nu_j(E_i)\\
    & = \sum\limits_{i=1}^{\infty} \nu(E_i)\\
\end{aligned}
\]
by exchanging the summation, which can be shown by taking the partial summations $\sum\limits_{i=1}^n \sum\limits_{j=1}^n a_jv_j(E_i)$.
\vspace{0.5em}

%%%%%%%%%%%%%%%%%%%%%%%%%%%%%%%%%%%%%%%%%%

\subsection*{Problem 2. Ex. 1.8. } 
Let $\{A_n\}$ be a sequence of events on a probability space $(\Omega, \mathcal{F}, P)$.\par
Define $\limsup_n A_n = \bigcap_{n=1}^{\infty}\bigcup_{i=n}^{\infty} A_i$ and $\liminf_n A_n = \bigcup_{n=1}^{\infty}\bigcap_{i=n}^{\infty} A_i$.\par
Show that $P(\liminf_n A_n) \leq \liminf_n P(A_n)$ and $\limsup_n P(A_n) \leq
P(\limsup_n A_n)$.\par
\vspace{0.5em}
\textbf{Sol.} \par
Notice
\[
\bigcap_{i=n}^{\infty} A_i\text{ increasing},\quad \text{and}\quad\bigcup_{i=n}^{\infty} A_i\text{ decreasing},
\]
we have
\[
P\left(\liminf_n A_n\right) = \limsup_n P\left(\bigcap_{i=n}^{\infty}A_i\right) \leq \limsup_n \inf_{k\geq n}P(A_k) = \liminf_n P(A_n) 
\]
and
\[
P\left(\limsup_n A_n\right) = \liminf_n P\left(\bigcup_{i=n}^{\infty}A_i\right) \geq \liminf_n \sup_{k\geq n}P(A_k) = \limsup_n P(A_n).
\]
\vspace{0.5em}

%%%%%%%%%%%%%%%%%%%%%%%%%%%%%%%%%%%%%%%%%%

\subsection*{Problem 3. Ex. 1.10. } 
Let $F(x_1, \cdots, x_k)$ be a c.d.f. on $\R^k$. Show that
\begin{enumerate}
    \item[(a)] $F(x_1, \cdots, x_{k-1}, x_k) \leq F(x_1, \cdots, x_{k-1}, x_k^\prime)$ if $x_k ≤\leq x_k^\prime$.\par
    \item [(b)] $\lim\limits_{ x_i\to -\infty} F(x_1, \cdots, x_k) = 0$ for any $1 \leq i \leq k$.
    \item [(c)] $F(x_1, \cdots, x_{k-1},\infty) = \lim\limits_{x_k\to\infty}F(x_1,\cdots, x_{k-1}, x_k)$ is a c.d.f. on $\R^{k−1}$.
\end{enumerate}
\par
\vspace{0.5em}
\textbf{Sol.} \par
There is a proability measure $P$ on $\R^k$ such that
\[F(x_1,\cdots,x_{k-1},x_k) = P((-\infty,x_1]\times \cdots (-\infty,x_{k-1}] \times (-\infty,x_k])\]\par
(a)We know
\[\begin{aligned}
    F(x_1,\cdots,x_{k-1},x_k) &= P((-\infty,x_1]\times \cdots (-\infty,x_{k-1}] \times (-\infty,x_k]) \\
    &\leq P((-\infty,x_1]\times \cdots (-\infty,x_{k-1}] \times (-\infty,x_k^\prime]) \\
    &= F(x_1,\cdots,x_{k-1},x_k^\prime).
\end{aligned}
\]\par
(b)
\[
\begin{aligned}
    \lim\limits_{ x_i\to -\infty} F(x_1, \cdots, x_k) &= \lim_{x_i\to -\infty}P((-\infty,x_1]\times \cdots (-\infty,x_{k-1}] \times (-\infty,x_k]) \\
    &= P\left(\bigcap_{x_i\to-\infty}(-\infty,x_1]\times \cdots (-\infty,x_{k-1}] \times (-\infty,x_k]\right) \\
    & = 0.
\end{aligned}
\]\par
(c) Define $\overline{P}$ to be the probability measure on $\R^{k-1}$ such that
\[\overline{P}((-\infty,x_1]\times \cdots (-\infty,x_{k-2}] \times (-\infty,x_{k-1}]) = P((-\infty,x_1]\times \cdots (-\infty,x_{k-1}] \times \R)\]
on $\R^{k-1}$, and whose existence can be secured by Carathoery's Theorem. Then
\[
\begin{aligned}
    F(x_1, \cdots, x_{k-1},\infty) &= \lim\limits_{x_k\to\infty} P((-\infty,x_1]\times \cdots (-\infty,x_{k-1}] \times (-\infty,x_k]) \\
    &=
     P\left((-\infty,x_1]\times \cdots (-\infty,x_{k-1}] \times \bigcup\limits_{x_k\to\infty}(-\infty,x_k]\right) \\
    &=
    P((-\infty,x_1]\times \cdots (-\infty,x_{k-1}] \times \R) \\
    &= \overline{P}((-\infty,x_1]\times \cdots (-\infty,x_{k-2}] \times (-\infty,x_{k-1}])
\end{aligned}
\]
and we are done.
\vspace{0.5em}

%%%%%%%%%%%%%%%%%%%%%%%%%%%%%%%%%%%%%%%%%%

\subsection*{Problem 4.} 

For a matrix $A \in \R^{n\times m}$, recall that the operator norm is the largest singular value: $\lVert A\rVert_{\text{op}} = s_{\max}(A)$.
Another useful norm is the Frobenius norm, defined by
\[\lVert A\rVert_F = \left(\sum\limits_{i=1}^n\sum\limits_{j=1}^m A_{ij}^2\right)^{1/2}.\]
Let $r = \text{rank}(A)$, and hence $1 \leq r \leq (n\wedge m)$. Write the nonzero singular values as $s_1(A) \geq \cdots \geq
s_r (A) > 0$.\par
(i) Prove that $\lVert A\rVert_F^2 = \sum\limits_{i=1}^r s_k(A)^2$.\par
(ii) Prove that $‖\lVert A‖\rVert_{\text{op}} \leq \lVert A\rVert_F \leq \sqrt{r} \lVert A\rVert_{op}$.\par
(iii) Show that $\lVert A\rVert_{\text{op}} = \lVert A\rVert_F$ if and only if $r = 1$. Characterize the nonzero eigenvalues of $AA^T$ and $A^TA$ when $\lVert A\rVert_F = \sqrt{r} \rVert A\rVert_{\text{op}}$.\par
\vspace{0.5em}
\textbf{Sol.} \par
(i) Consider the SVD of $A = QSV$ where $Q \in \R^{n\times n}, V\in \R^{m\times m}$ are orthogonal matrices. Then
\[
\begin{aligned}
\sum\limits_{i=1}^n\sum\limits_{j=1}^m A_{ij}^2 &= \sum\limits_{i=1}^n |A_i|^2 \\
&= \sum\limits_{i=1}^n |(Q_iSV)|^2 \\
&= \sum\limits_{i=1}^n \sum\limits_{k=1}^nq_{ik}^2s_k^2 =  \sum\limits_{k=1}^n s_k^2 = \sum\limits_{k=1}^r s_k(A)^2
\end{aligned}
\]
since $Q,V$ orthogonal and denote $s_k := s_k(A)$ for $1\leq k \leq r$ and $s_k:= 0$ for $r<k\leq n$.\par
(ii) Notice
\[\lVert A\rVert_{\text{op}}^2 = s_{\max}^2(A) \leq \sum\limits_{i=1}^r   s_i^2(A) = \lVert A\rVert_{F}^2 \leq rs_{\max}^2(A) = r\lVert A\rVert_{\text{op}}^2\]
and we are done.\par
(iii) If $r = 1$, then obviously
\[
    \lVert A\rVert_{\text{op}}^2 = s_1(A)^2 = \lVert A\rVert_{\text{F}}^2.
\]
For the other direction, we know in this case we have
\[
\lVert A\rVert_{\text{op}}^2 = s_{\max}^2(A) = \sum\limits_{i=1}^r s_k^2(A) = \lVert A\rVert_{F}^2
\]
that means $s_k(A) = 0$ for all $k\geq 2$ if $r\geq 2$, which is a contradiction and hence $r = 1$.\par
If $\lVert A\rVert_{F} = \sqrt{r}\lVert A\rVert_{\text{op}}^2$, we know
\[
rs_1(A)^2 = \sum\limits_{i=1}^r   s_i^2(A)
\]
and hence $s_1(A) = s_2(A) = \cdots = s_r(A) = \lVert A\rVert_{\text{op}}$ and since
\[
AA^T = QS^2Q^T,\quad A^TA = V^TS^2V
\]
and we know the eigenvalues of $AA^T$ and $A^TA$ will al equal to $s_1(A)^2 = \lVert A\rVert_{\text{op}}^2$.
\vspace{0.5em}

%%%%%%%%%%%%%%%%%%%%%%%%%%%%%%%%%%%%%%%%%%

\subsection*{Problem 5.} 

Recall the example from the lecture: let $U \sim \text{Unif}([0, 2\pi])$, and define $X = \cos U$ and $Y = \sin U$ . We want to describe their joint distribution $\mu_{X,Y}$ on $\R^2$.\par
(1) Define $f (u) = (\cos u, \sin u)$ for $u \in [0, 2\pi]$. Prove that $f$ is measurable as a map:\[f:([0, 2\pi], \B([0, 2\pi])) \to (\R^2, \B(\R^2)).\] You may cite an appropriate theorem. Use the definition of an induced measure to express
$\mu_{X,Y}(B)$ for $B \in \B(\R^2)$.\par
(2) Alternatively, construct the distribution as a measure on the unit circle. Let
\[S = \{(x, y) \in \R^2 : x^2 + y^2 = 1\},\quad  A_{u,v} = \{(\cos z, \sin z) : z \in (u, v)\}, u < v.\]
Let $\B(S)$ be the $\sigma$-algebra generated by all sets $A_{u,v}$ , where $u, v \in \R$ and $u < v$.\par
(i) Prove that there is a unique probability measure $\mu_0$ on $(S, \B(S))$ such that
\[\mu_0(A_{u,v} ) = \dfrac{v - u}{2\pi}\quad\text{ for every } 0 < u < v < 2\pi.\]
You may cite an appropriate measure extension or uniqueness theorem.\par
(ii) For every $B \in B(\R^2)$, show that $B \cap S \in \B(S)$. Then show that
\[
\mu(B) := \mu_0(B \cap S), B \in \B(\R^2),
\]
defines a probability measure on $\R^2$.\par
(iii) Show that $\mu = \mu_{X,Y}$, where $\mu_{X,Y}$ is the measure from part (1).
\par
\vspace{0.5em}
\textbf{Sol.} \par
(1) We only need to prove that $f^{-1}(A)$ is measurable for any $A\in\mathcal{A}$ where $\sigma(\mathcal{A}) = \B(\R^2)$ for a basis $\mathcal{A}$. This conclusion maybe found on any measure theory based probability theory or a measure theory book, so I will skip the proof of that.\par
Let $\mathcal{A}$ to be the set of closed balls in $\R^2$, and if $B\in \mathcal{A}$, we know $B\cap T$ where $T$ is the unit cycle on $\R^2$ is always a close arc, which means $f^{-1}(B)$ is always a closed interval or a closed interval with a singleton $\{2\pi\}$, both belong to $\B([0,2\pi])$ and we are done.(The detailed proof my refer 2(ii))\par
Notice
\[\mu_{X,Y}(B) = P((X,Y) \in B) = P(f(U) \in B) = P(f^{-1}(B))\]
where $P$ is the probability measure corresponding to $\text{Unif}([0,2\pi])$.\par
2. (i) For any closed interval $[a,b] \subset [0,2\pi]$, define
\[\rho([a,b]) = \dfrac{b-a}{2\pi}\]
and define
\[
\mu^*(B) = \inf\left\{\sum\limits_{i=1}^{\infty} \rho(I_i), \text{ where }I_i\text{ are closed intervals in }[0,2\pi], B\subset \bigcup_{i=1}^{\infty} I_i\right\}
\]
which is an out measure on $[0,2\pi]$, and it is easy to check for any $E\subset [0,2\pi]$ and $I$ and closed interval in $[0,2\pi]$,
\[
\mu^*(E) = \mu^*(E\cap I) + \mu^*(E\cap I^c)
\]
TO check this, we only need to prove LHS greater than RHS, consider $\{E_i\}_{i=1}^{\infty}$ an closed intervals cover of $E$, then we know $E_i \cap I$ is an closed intervals and $E_i \cap (\mathring{I})^c$ are disjoint union of two closed intervals, with $\mu^*(E_i \cap I) + \mu^*(E_i \cap (\mathring{I})^c) = \mu^*(E_i)$ and it will induced the required conclusion. So any closed interval are $\mu^*$-measurable, which means $\B([0,2\pi])$ are $\mu^*$-measurable and $\mu^*$ is a measure on $[0,2\pi]$ by Caratheodory's Theorem.\par
Now consider the induced measure $\mu_0$ on $S$, there exists a natural inclusion $\iota:S\to[0,2\pi]$ and we define $\mu_0(E): = \mu^*(\iota(E))$ for $E\in \B(S)$ where we may prove $\iota(\B(E)) \subset \B([0,2\pi])$ by $\pi$-$\lambda$ theorem, that is we know the closed arcs on $S$ is a $\pi$-system and the set $E$ with $\iota(E) \in \B([0,2\pi])$ will form a $\lambda$-system, so $\mu_0$ is well defined on $\B(S)$ and if there is another measure $\nu$ satisfies the equality, then we use the $\pi$-$\lambda$ Theorem again by considering
\[\mathcal{M} = \{E\in\B(S), \mu_0(E) = \nu(E)\}\]
and we know the closed arcs is a subset of $\mathcal{M}$ and a $\pi$-system at the same time. We would like to show $\mathcal{M}$ to be a $\lambda$-system, $S\in \mathcal{M}$ and if $A\subset B, A,B\in\mathcal{M}$, then $\mu_0(B-A) = \mu_0(B)- \mu_0(A) = \nu(B)-\nu(A) = \nu(B-A)$ obviously. For $A_i \uparrow A$, then we know
\[\mu_0(A) = \limsup \mu_0(A_i) = \limsup \nu(A_i) = \nu(A)\]
and hence $\mathcal{M}$ is a $\lambda$-system. Therefore, $\mu_0 = \nu$ for any set in $\B(S)$ are we are done.\par
(ii) We still use the same notation $\mathcal{A}$ in (1) to denote all the closed balls in $\R^2$ and we know $B\cap S$ is a closed arc and hence in $\B(S)$ for any $B\in\mathcal{A}$, so basically we proved the claim in (1) here again. Now consider $\mathcal{P} = \{B\in\B(\R^2), B\cap S \in\B(S)\}$ and we know $\mathcal{A}\subset\mathcal{P}$ with the fronter a $\pi$-system, so we prove $\mathcal{P}$ to be a $\lambda$ system, and obviously $\R^2 \in \mathcal{P}$. For $A\subset B, A,B\in\B(\R^2)$, we know $(B-A)\cap S = B\cap S-A\cap S$ and for $A_i \uparrow A_i$, $A\cap S = \bigcup_{i=1}^{\infty} (A_i \cap S)$ and hence $\mathcal{P}$ is a $\lambda$-system. Therefore, $\B(\R^2) = \sigma(\mathcal{A}) \subset \mathcal{P}$ and we are done.\par
(iii) For any $B\in\B(\R^2)$, we know
\[
\mu(B) = \mu_0(B\cap S) = \mu^*(\iota(B\cap S)),\quad \mu_{X,Y}(B) = P(f^{-1}(B))
\]
and we know $\iota(B \cap S) = f^{-1}(B)$ $P$-almost surely and $\mu^*$-almost surely, so we only need to check $\mu^* = P$ on $\B([0,2\pi])$, with the same techniques in (i), since
\[\mu^*([u,v]) = \dfrac{u-v}{2\pi} = P([u,v])\]
for any closed interval $[u,v]$ in $[0,2\pi]$, we know they are the same and we are done.
\vspace{0.5em}





\end{document}
